%% =========================================================
%% Chapter: The conditional skew estimate
%% Source: OpenAI, "A two-dimensional area law from a global spectral gap"
%%         (September 24, 2026), Lemma 5.3, 04-conditional.tex lines 487-600
%% =========================================================

\chapter{The Conditional Skew Estimate}
\label{ch:conditional_skew}

This chapter formalizes the conditional skew estimate
of~\cite[Lemma~5.3]{OpenAI2026AreaLaw}: a local positive expectation
conjugated by marginal real powers departs from the nonnegative real axis by
at most a fractional power of a conditional mutual information. Throughout,
$\theta$ is a unit vector on finite-dimensional systems $P,x,U,F$ with
$P=P_0P_1$, $Y=xU$, $\rho_D$ denotes a marginal of $\theta$, powers
$\rho^{[z]}$ vanish on the kernel, and $0\le h\le I$ acts on $P_0x$. Put
$d_h=(\dim P_0)(\dim x)$, $\ell_h=\log(ed_h)$, $\eta=I(x:F|U)_\theta$ and
\[
  f(z)=\langle\theta,\rho_P^{[z]}\rho_Y^{[-z]}h\rho_P^{[-z]}\rho_Y^{[z]}\theta\rangle .
\]

\section{The scalar function}

\begin{lemma}[Elementary properties]
    \label{lem:skew_function}
    \lean{Entropy.ConditionalSkew.differentiable_skewFun,
        Entropy.ConditionalSkew.skewFun_zero,
        Entropy.ConditionalSkew.skewFun_neg_conj,
        Entropy.ConditionalSkew.norm_skewFun_imag_le,
        Entropy.ConditionalSkew.norm_skewFun_le_of_strip}
    \leanok
    The function $f$ is entire, $f(0)=\langle\theta,h\theta\rangle\ge0$ and
    $f(-\bar z)=\overline{f(z)}$. On the imaginary axis $|f|\le1$, and
    $|f(z)|\le d_h^2$ for $|\operatorname{Re}z|\le\frac12$.
\end{lemma}

\begin{proof}\leanok
    Each matrix entry of a support power is an entire function of the exponent.
    On the imaginary axis $f(iy)$ is an expectation of $h$ in a unit vector.
    For the strip bound, expand $h$ in matrix units on $P_0x$; each of the at
    most $d_h^2$ terms is an inner product of two vectors of the form
    $\rho^{[z]}c\,\rho^{[-z]}\theta$ with $\lVert c\rVert\le1$, and
    $\lVert\rho^{[z]}c\,\rho^{[1/2-z]}\rVert_2\le\lVert c\rVert$ for
    $0\le\operatorname{Re}z\le\frac12$. The reflection treats negative real
    parts.
\end{proof}

\begin{theorem}[Departure from the positive axis]
    \label{thm:skew_departure}
    \lean{Complex.norm_sub_re_le_of_phase_bounds}
    \leanok
    Let $g$ be entire with $g(-\bar z)=\overline{g(z)}$, $g(0)=p\ge0$,
    $|g|\le1$ on the imaginary axis and $|g|\le m$ on
    $|\operatorname{Re}z|\le\frac14$, where $1\le m$, $1\le\ell$ and
    $\log m\le2\ell$. Suppose that for some $R\ge0$ and every real $u$,
    \[
      |g(iu)-p|\le4\sqrt p\,|\sinh\pi u|\,R+4\sinh^2(\pi u)R^2 .
    \]
    Then for real $|t|\le1/(16\ell)$,
    \[
      |g(t)|-\operatorname{Re}g(t)\le t^2\bigl(8\pi^2R^2
        +12288\,e^{\pi^2+1}\ell^2\sqrt{\min(1,R)}\bigr).
    \]
\end{theorem}

\begin{proof}\leanok
    Three lines on the two half-strips bounds $g$ by a universal constant on
    $|\operatorname{Re}z|\le1/(8\ell)$. The phase hypothesis gives
    $g'(0)\in i\mathbb R$ and $|g'(0)|\le4\pi\sqrt p\,R$. The function
    $(g(z)-p)e^{z^2}$ is bounded by $C\min(1,R)$ on the imaginary axis and by a
    constant on the outer lines, so three lines and a Cauchy estimate bound the
    second-order Taylor remainder of $g$ at zero. Finally
    $|p+itb|-p\le t^2b^2/p$ for real $b$ when $p>0$, and $g'(0)=0$ when $p=0$.
\end{proof}

\section{The imaginary axis}

\begin{lemma}[Schmidt mirroring]
    \label{lem:schmidt_mirroring}
    \lean{Entropy.ConditionalSkew.cfcC_kronecker_one_mulVec}
    \leanok
    For a vector $\theta$ on $A\otimes B$ and any function $g$ on the real
    line, $(g(\rho_A)\otimes I)\theta=(I\otimes g(\rho_B))\theta$.
\end{lemma}

\begin{proof}\leanok
    Write $\theta$ as a coefficient matrix $C$, so $\rho_A=CC^*$ and
    $\rho_B=(C^*C)^T$. The operators act by $C\mapsto g(CC^*)C$ and
    $C\mapsto C\,g(C^*C)$, which agree since $CC^*\cdot C=C\cdot C^*C$.
\end{proof}

\begin{lemma}[Imaginary-axis comparison]
    \label{lem:skew_imaginary_axis}
    \lean{Entropy.ConditionalSkew.norm_skewFun_imag_sub_le}
    \leanok
    \uses{lem:skew_function, lem:schmidt_mirroring, thm:marginal_phase}
    With $p=\langle\theta,h\theta\rangle$ and $\mathcal R_\eta$ the rate of
    Theorem~\ref{thm:marginal_phase} for $d=\dim x$, every real $u$ satisfies
    \[
      |f(iu)-p|\le4\sqrt p\,|\sinh\pi u|\,\mathcal R_\eta
        +4\sinh^2(\pi u)\mathcal R_\eta^2 .
    \]
\end{lemma}

\begin{proof}\leanok
    The kernel parts of the support powers annihilate $\theta$, and Schmidt
    mirroring across $P$ and $YF$ gives $f(iu)=\langle w_u,hw_u\rangle$ with
    $w_u=\widehat\rho_Y^{\,iu}\widehat\rho_{YF}^{\,-iu}\theta$. The vector
    $w_u^0=\widehat\rho_U^{\,iu}\widehat\rho_{UF}^{\,-iu}\theta$ is a unitary
    on $UF$ applied to $\theta$, so $\langle w_u^0,hw_u^0\rangle=p$.
    Multiplying $w_u-w_u^0$ by $\widehat\rho_P^{\,iu}$ and
    $\widehat\rho_U^{\,-iu}$ and mirroring again gives
    $\lVert w_u-w_u^0\rVert=\lVert(\widehat\rho_U^{\,-iu}\widehat\rho_Y^{\,iu}
    -\widehat\rho_{UF}^{\,-iu}\widehat\rho_{YF}^{\,iu})\theta\rVert$, which
    Theorem~\ref{thm:marginal_phase} bounds by
    $2|\sinh\pi u|\,\mathcal R_\eta$. For $0\le h\le I$,
    $|\langle w,hw\rangle-\langle w^0,hw^0\rangle|\le
    2\sqrt{\langle w^0,hw^0\rangle}\,\lVert w-w^0\rVert+\lVert w-w^0\rVert^2$.
\end{proof}

\section{The estimate}

\begin{lemma}[Size of the rate]
    \label{lem:skew_rate}
    \lean{Entropy.ConditionalSkew.phaseRate_le,
        Entropy.ConditionalSkew.condMutualInfo_nonneg,
        Entropy.ConditionalSkew.condMutualInfo_le_two_log}
    \leanok
    For a density matrix on $xUF$, $0\le I(x:F|U)\le2\log\dim x$. If
    $1\le D$, $\log(eD)\le\ell$ and $\eta\le2\ell$, then the rate of
    Theorem~\ref{thm:marginal_phase} satisfies
    $\mathcal R_\eta\le3\ell\min(1,\eta)^{1/4}$.
\end{lemma}

\begin{proof}\leanok
    Nonnegativity is strong subadditivity, and the upper bound writes
    $I(x:F|U)=S(x|U)-S(x|UF)$ with both conditional entropies bounded in
    absolute value by $\log\dim x$. For $0<\eta\le1$ use
    $-\eta\log\eta\le2\sqrt\eta$; for $\eta>1$ the rate is at most
    $1+\sqrt{2}\,\ell$.
\end{proof}

\begin{theorem}[Conditional skew estimate]
    \label{thm:conditional_skew}
    \lean{Entropy.ConditionalSkew.conditionalSkew_le}
    \leanok
    \uses{lem:skew_function, thm:skew_departure, lem:skew_imaginary_axis,
        lem:skew_rate}
    There is a universal constant $C$ such that for $0<a\ell_h\le\frac18$ and
    $t=a/2$,
    \[
      |f(t)|-\operatorname{Re}f(t)\le Ca^2\ell_h^4\eta^{1/8}.
    \]
    In particular the left side is zero if $\eta=0$. One may take
    $C=18\pi^2+6144\,e^{\pi^2+1}$.
\end{theorem}

\begin{proof}\leanok
    Apply Theorem~\ref{thm:skew_departure} to $f$ with $m=d_h^2$,
    $\ell=\ell_h$, $R=\mathcal R_\eta$ and $p=\langle\theta,h\theta\rangle$,
    using Lemmas~\ref{lem:skew_function} and~\ref{lem:skew_imaginary_axis}.
    Since $\dim x\le d_h$ and $\eta\le2\log\dim x\le2\ell_h$,
    Lemma~\ref{lem:skew_rate} gives
    $\mathcal R_\eta^2\le9\ell_h^2\eta^{1/8}$ and
    $\sqrt{\min(1,\mathcal R_\eta)}\le2\ell_h\eta^{1/8}$.
\end{proof}
